% META: {"id":"1SPE-DERLOCAL-EX-001","chapitre":"1SPE-DERIVATION-LOCAL","type_objet":"exercice","capacites":["1SPE-DERIVATION-LOCAL-C1"],"methodes":["M1"],"parcours":1,"competences":["calculer"],"duree_min":5,"mode_creation":"ex_nihilo","sources_inspiration":[],"parametres_sympy":{"a":1,"b":4},"fichier_tex":"chapitres/1SPE-DERIVATION-LOCAL/exercices/1SPE-DERLOCAL-EX-001.tex","corrige_tex":"chapitres/1SPE-DERIVATION-LOCAL/corriges/1SPE-DERLOCAL-CO-001.tex","status":"approved","origin":"generated","status_history":["generated","audited","accepted","approved"],"release_acceptance":"RELEASE_OWNER_BATCH_ACCEPTANCE","acceptance_closure_digest":"sha256:9b3ccf9a81c5520fbb7e03b7b2d3e7bf2057a02834b6d908bf3be5f49f3b553f","acceptance_receipt_digest":"sha256:4fad86f0282e0fa4e0419cca1ad6379ae7af5a280c04a4a90880f90a1c044242"}
% BEGIN-VERIFY
% from sympy import Rational
% assert Rational((4**2+1)-(1**2+1), 4-1) == 5
% END-VERIFY
\begin{exercice}{1SPE-DERLOCAL-EX-001}{1}{5}
Soit $f(x)=x^2+1$. Calculer le taux de variation de $f$ entre $1$ et $4$.
\end{exercice}
